\FloatBarrier
\section{本章小结}
\label{sec:direct-methods-summary}

本章回答了章首的问题：对非奇异方阵，怎样不求逆、在有限步内把精确解算出来。共同的做法是用消元把方阵化为三角方程，并把分解与右端的处理分开，使同一次分解可以服务于多个右端。

消元是这些方法的起点。初等行变换不改变解，把增广矩阵化为阶梯形就能读出主元、自由变量与秩；对$[\bm A\mid\bm I_n]$消元可以求逆，但只解一个方程时不必先求逆。

有了分解，求解就化为前代与回代：前代自上而下求出$\bm L\bm y=\bm P\bm b$的各分量，回代自下而上求出$\bm U\bm x=\bm y$的各分量。一般非奇异方阵用带部分选主元的LU分解$\bm P\bm A=\bm L\bm U$，消元逐列进行，处理完前$d-1$列后结束，乘子按位置存入$\bm L$，后续换行时已存乘子要随行交换；部分选主元对任何非奇异矩阵都能进行，不选主元则要求各顺序主子式非零。

利用对称性可以减少一半的工作。对称正定矩阵用Cholesky分解$\bm A=\bm L\bm L^\top$，每一列的开方量是一个Schur补，正定性恰好保证它为正；分解存在且唯一当且仅当矩阵对称正定，运算量是LU的一半。对称不定矩阵用分块LDL分解$\bm P\bm A\bm P^\top=\bm L\bm D\bm L^\top$，负主元保留在$\bm D$中，零对角元时改用二阶块，使可逆的矩阵总能继续消元。

低秩修改让已有分解得以复用。低秩更新把修改的影响集中到$k$维辅助量上，只需旧因子上的$k+1$次求解和一个$k$阶方程；更新后的矩阵可逆当且仅当辅助矩阵$\bm K$可逆，Woodbury恒等式和Sherman--Morrison公式是这一等价关系的矩阵形式。这种节省的前提是已有分解且修改秩很小。

无论用哪种方法，结果都要代回原方程检验$\bm A\bm x=\bm b$。这一检验说明算出的$\bm x$符合目标，却不说明它离真解有多近；主元增长在本章只作了算法层面的说明，结果的可信程度留给第~\ref{chap:solution-error-regularization}章。不相容或非方阵的问题要改求最小二乘解，这是下一章的内容。
